% ------------------------------------------------------------------------------
\section{Descargar ConcretoLab}
\label{sec:descarga}
% ------------------------------------------------------------------------------

ConcretoLab se distribuye desde el repositorio oficial del proyecto. El instalador es un único archivo,
\ruta{\archivoinstalador}, de unos 115~MB, que contiene el programa y todo lo que necesita para funcionar.

\subsection{Enlace de descarga}

\paso{1}{Abra el enlace de descarga}

Haga clic en el siguiente enlace o cópielo en la barra de direcciones de su navegador (Microsoft Edge, Google Chrome o
Firefox). La descarga comienza de inmediato y siempre corresponde a la versión más reciente:

\begin{center}
	\href{\enlacedescarga}{\ruta{\enlacedescarga}}
\end{center}

Desde un celular o una tableta puede escanear este código, aunque el programa solo se instala en una computadora con
Windows:

\begin{figure}[H]
	\centering
	\qrcode[height=3.2cm]{\enlacedescarga}
	\caption{Código QR del enlace de descarga.}
	\label{fig:qr-descarga}
\end{figure}

\paso{2}{Espere a que termine la descarga}

El navegador guarda el archivo en la carpeta \ruta{Descargas}. Según la velocidad de su conexión, la descarga demora de
uno a varios minutos. Cuando termine, el archivo aparece con el ícono de ConcretoLab:

\begin{figure}[H]
	\centering
	\begin{tabular}{c}
		\includegraphics[width=1.6cm]{guia/00-icono-instalador} \\
		\ruta{\archivoinstalador}
	\end{tabular}
	\caption{El instalador descargado, en la carpeta Descargas.}
	\label{fig:icono}
\end{figure}

\paso{3}{Si el navegador muestra un aviso}

Algunos navegadores advierten que el archivo «no se descarga habitualmente» o piden confirmar la descarga. Ocurre con
los programas que todavía no tienen firma digital de un fabricante reconocido. En Microsoft Edge abra el menú de
descargas, pulse los tres puntos junto al archivo y elija \boton{Conservar}; si pregunta de nuevo, pulse
\boton{Mostrar más} y luego \boton{Conservar de todos modos}. En Google Chrome pulse \boton{Conservar}.

\subsection{Versiones anteriores y notas de cada versión}

En la página de versiones del repositorio están todas las versiones publicadas, con la fecha, los cambios de cada una y
la huella SHA-256 de su instalador:

\begin{center}
	\href{\enlaceversiones}{\ruta{\enlaceversiones}}
\end{center}

Para comprobar que el archivo descargado está completo y no fue modificado, siga el procedimiento del
anexo~\ref{sec:huella}. Es opcional.

\subsection{Si no tiene internet en la computadora de trabajo}

Descargue el instalador en otra computadora y cópielo en una memoria USB. El programa no necesita internet para
instalarse ni para funcionar.
